\section{Conclusion}

In this work, we presented OmniRetrieval, a framework for retrieval over structurally heterogeneous knowledge sources that, given a natural-language question, engages each relevant source through its own native query language and consolidates the executor outputs via a cross-source evidence selection step, rather than collapsing the sources into a shared representation. Evaluation on a benchmark spanning 13 datasets and 309 knowledge bases over unstructured corpora, relational databases, RDF graphs, and labeled property graphs shows that OmniRetrieval consistently outperforms relevant baselines. Our analyses further indicate that broad exploration at the source-selection step, with the final commitment deferred to a selector that rests on retrieved evidence, is what lets OmniRetrieval scale gracefully. These findings position OmniRetrieval as a step toward a general-purpose universal layer, one that preserves the structural affordances that make each source valuable while exposing a single natural-language interface to the user.

\section*{Limitations}

OmniRetrieval is a general framework, and our specific instantiation in this work opens a couple of directions for further exploration. First, while our instantiation of the cross-source evidence selection step already performs reliably, it would be valuable to further strengthen this in future work, for example, through supervised fine-tuning on labeled cross-source selections or reinforcement learning that uses downstream answer quality as the reward signal. Second, our setup uses a single shared LLM across source selection, native query formulation, and evidence selection, and operator-specific specialization is another avenue worth investigating.

\section*{Ethical Considerations}

Our work builds on publicly available datasets and standard LLMs, and we do not foresee ethical concerns beyond those inherited from these underlying components. In particular, since the retrieved evidence and the generated source-native queries draw on the connected knowledge bases and the internalized knowledge within the LLMs, any private, harmful, or biased content present in these resources may carry over to the outputs, and we recommend applying standard safeguards and filtering over both the connected sources and the generated queries to ensure safe and responsible deployment.
